\chapter{The Sell Side}\label{ch:m1:the-sell-side}

At 11:40 a pension fund decides to sell two million shares of a stock that
trades ten million a day. Worked patiently, the order would take two days and
the price might fall all the while. The fund's trader phones a bank instead.
Forty seconds later she has a price for the whole block, a dollar and a half
below the screen, and says ``done''. The bank now owns a hundred million
dollars of stock it never wanted. This chapter is about the firms that say
yes to that call, how they decide the price, and why they are organised the
way they are.

\section{The markets division}

\begin{definition}[Sell side and broker-dealer]\label{def:m1:the-sell-side:sell-side}
The \emph{sell side}\index{sell side} is the set of firms that sell trading
services to investors: execution, prices, financing, research and new
issues. Its legal vehicle is the \emph{broker-dealer}\index{broker-dealer}, a
firm licensed both to execute clients' orders as agent (broker) and to trade
with them as principal (dealer).
\end{definition}

The largest \omterm{def:m1:the-sell-side:sell-side}{broker-dealers} sit inside banks. A bank's \emph{markets
division} is organised around the client, not around a strategy:

\begin{itemize}
\item \textbf{Sales} owns the client relationship and brings in the enquiry.
\item \textbf{Trading} makes the price and manages the resulting risk.
\item \textbf{Structuring} designs products that package several risks for a
  client --- the subject of One Quant Book 5.
\item \textbf{Research} publishes analysis that gives clients a reason to
  call.
\item \textbf{Quants and technology} build the pricing, risk and execution
  systems the other four run on.
\end{itemize}

\begin{definition}[Sales-trader]\label{def:m1:the-sell-side:sales-trader}
A \emph{sales-trader}\index{sales-trader} takes clients' orders and works
them in the market on their behalf, as agent, advising on timing and
tactics; a \emph{trader}, by contrast, commits the bank's capital.
\end{definition}

\begin{omfigure}
\begin{tikzpicture}[font=\small,
  box/.style={draw=omDef,rounded corners=2pt,minimum width=2.7cm,minimum height=0.85cm,align=center,fill=omDef!4},
  ext/.style={draw=omExo,rounded corners=2pt,minimum width=2.7cm,minimum height=0.85cm,align=center,fill=omExo!6},
  flow/.style={-{Stealth[length=2mm]},thick,omDef}]
\node[ext] (cl) at (0,0) {Client};
\node[box] (sa) at (4.2,1.5) {Sales};
\node[box] (st) at (4.2,0) {Sales-trading\\(agent)};
\node[box] (re) at (4.2,-1.5) {Research};
\node[box] (tr) at (8.6,0.75) {Trading\\(principal)};
\node[box] (su) at (8.6,-0.75) {Structuring};
\node[ext] (mk) at (13,0.75) {Markets};
\node[box] (qt) at (8.6,-2.4) {Quants and technology};
\draw[flow] (cl) -- node[above,sloped]{enquiry} (sa);
\draw[flow] (cl) -- node[above]{orders} (st);
\draw[flow] (re) -- node[below,sloped]{ideas} (cl);
\draw[flow] (sa) -- (tr);
\draw[flow] (st) -- (tr);
\draw[flow] (tr) -- node[above]{hedges} (mk);
\draw[flow] (su) -- (tr);
\draw[dashed,omExo] (qt) -- (su);
\end{tikzpicture}

\omcaption{A markets division seen from the client. Everything in the middle
column exists to bring business to the right-hand column, which carries the
risk.}
\end{omfigure}

\begin{definition}[Flow trading and franchise]\label{def:m1:the-sell-side:flow}
\emph{Flow trading}\index{flow trading} is market making to clients: the desk
quotes on request, takes the other side, and manages the inventory that
results. The desk's \emph{franchise}\index{franchise} is the stream of client
enquiry it sees; its value is both the spread earned on that flow and the
information the flow carries about who wants to do what.
\end{definition}

\begin{remark}[A bank desk is not a proprietary firm]\label{rem:m1:the-sell-side:notprop}
An electronic \omterm{def:m1:what-a-trading-firm-does:market-maker}{market maker} (\cref{ch:m1:what-a-trading-firm-does}) competes
on price and speed for anonymous orders. A bank desk competes for
\emph{relationships}: the client who is shown a good price in a difficult
block also gives the bank its financing, its derivatives and its next bond
issue. A trade that loses money in isolation can be good business; deciding
when is the head of desk's job (One Quant Book 16).
\end{remark}

\section{Pricing a block}

\begin{definition}[Block trade and risk bid]\label{def:m1:the-sell-side:block}
A \emph{block trade}\index{block trade} is a transaction much larger than the
market's displayed size, negotiated away from the order book. In a
\emph{risk bid} the dealer buys the whole block as principal at a discount
to the market price; the discount pays it for the cost and the risk of
selling the shares afterwards.
\end{definition}

The dealer's problem has two parts. Selling a large quantity pushes the price
down: an \emph{impact} cost. And while it sells, the price moves for reasons
of its own: a \emph{risk}.

\begin{method}[The square-root rule of thumb for impact]\label{met:m1:the-sell-side:sqrt}
The average cost of trading $Q$ shares of a stock with daily volatility
$\sigma$ and daily volume $V$, as a fraction of their value, is close to
\[
  I(Q) \;=\; \kappa\,\sigma\sqrt{Q/V},
\]
with $\kappa$ of order one and, within reasonable limits, little dependence
on how long the execution takes. The empirical law and its theory belong to
One Quant Book 10; here it is a tool, with $\kappa = 0.7$.
\end{method}

\begin{proposition}[Risk of a linear unwind]\label{prop:m1:the-sell-side:tover3}
A dealer holds a block worth $N$ and sells it at a constant rate over $T$
days, while the price follows a driftless random walk with daily volatility
$\sigma$. Ignoring impact, its profit relative to selling everything at the
initial price has mean zero and standard deviation
\[
  N\,\sigma\sqrt{T/3}.
\]
\end{proposition}
\begin{proof}
Let $W_t$ be the price return since purchase, a Brownian motion with
variance $\sigma^2 t$. Selling at rate $1/T$, the dealer realises the average
return $\bar W = \frac1T\int_0^T W_t\dd t$. Then $\E[\bar W]=0$ and
\[
\Var(\bar W) = \frac{1}{T^2}\int_0^T\!\!\int_0^T \sigma^2 \min(s,t)\dd s\dd t
 = \frac{\sigma^2}{T^2}\cdot\frac{T^3}{3} = \frac{\sigma^2 T}{3}.
\qedhere
\]
\end{proof}

The factor $1/3$ is worth remembering: a position that shrinks linearly to
zero over $T$ days carries the risk of a constant position held for $T/3$
days.

\begin{proposition}[Break-even discount]\label{prop:m1:the-sell-side:discount}
With participation rate $\pi$ (the fraction of daily volume the dealer is
willing to be), the unwind lasts $T = Q/(\pi V)$ days, and the discount at
which the dealer loses money with probability $1-\Phi(z)$ only is
\[
  d_z \;=\; \kappa\,\sigma\sqrt{Q/V} \;+\; z\,\sigma\sqrt{\frac{Q}{3\pi V}} .
\]
\end{proposition}
\begin{proof}
The dealer's profit per unit of value is $d - I(Q) + \bar W$, normal with
mean $d - I(Q)$ and standard deviation $\sigma\sqrt{T/3}$ by
\cref{prop:m1:the-sell-side:tover3}. It is negative with probability
$\Phi\bigl(-(d-I)/(\sigma\sqrt{T/3})\bigr)$; set this equal to $\Phi(-z)$.
\end{proof}

\begin{example}[The pension fund's block]\label{ex:m1:the-sell-side:block}
$Q = 2$ million shares at \$50, $V = 10$ million, $\sigma = 2\%$ a day,
$\pi = 10\%$. Then $T = 2$ days, $I = 0.7 \times 0.02 \times \sqrt{0.2} =
0.63\%$, and the risk term at 95\% ($z = 1.645$) is $1.645 \times 0.02
\times \sqrt{2/3} = 2.69\%$. The break-even discount is $3.31\%$, or \$1.66 a
share: the dealer who bids \$48.34 expects to make $2.69\%$ of \$100
million, \$2.7 million, and loses money one time in twenty. Both terms grow
like $\sqrt{Q}$: \emph{a block four times larger costs twice as much per
share.}
\end{example}

\begin{omfigure}
\begin{tikzpicture}
\begin{axis}[width=11cm,height=5.4cm,
  xlabel={block size (\% of daily volume)},ylabel={discount (\% of value)},
  tick label style={font=\small},label style={font=\small},
  xmin=0,xmax=50,ymin=0,ymax=8.5,grid=major,grid style={gray!25},
  legend columns=4,legend style={font=\footnotesize,at={(0.5,-0.3)},anchor=north,draw=none,fill=none,/tikz/every even column/.append style={column sep=8pt}},
  table/col sep=comma]
\addplot[omThm,very thick] table[x=pct_adv,y=vol3]{figdata/markets-1/02-the-sell-side/discount_vs_size.csv};
\addplot[omDef,very thick] table[x=pct_adv,y=vol2]{figdata/markets-1/02-the-sell-side/discount_vs_size.csv};
\addplot[omMeth,very thick] table[x=pct_adv,y=vol1]{figdata/markets-1/02-the-sell-side/discount_vs_size.csv};
\addplot[omDef,dashed] table[x=pct_adv,y=impact2]{figdata/markets-1/02-the-sell-side/discount_vs_size.csv};
\legend{$\sigma=3\%$,$\sigma=2\%$,$\sigma=1\%$,impact only ($\sigma=2\%$)}
\end{axis}
\end{tikzpicture}

\omcaption{The 95\% break-even discount of a risk bid, at 10\% participation.
Most of the discount pays for risk, not for impact. Data: computed by the
chapter's script from \cref{prop:m1:the-sell-side:discount}.}
\end{omfigure}

\begin{remark}[What the model leaves out]\label{rem:m1:the-sell-side:limits}
Three things, all of which a desk prices by hand. \emph{Information:} a
client who sells because it knows something is worth a wider discount ---
\omterm{def:m1:what-a-trading-firm-does:adverse-selection}{adverse selection} again. \emph{Hedging:} the dealer can sell index futures
within seconds and keep only the stock-specific risk: a hedge that removes
half of the variance shrinks $\sigma$ by 29\%. \emph{\omterm{def:m1:the-sell-side:flow}{Franchise}:} the desk may
bid through its break-even to win the client's other business.
\end{remark}

\section{Financing, fees and new issues}

Making prices is one of three businesses in the division.

\begin{definition}[Prime services]\label{def:m1:the-sell-side:prime}
\emph{Prime services}\index{prime services} is the bank business that
finances \omterm{def:m1:what-a-trading-firm-does:hedge-fund}{hedge funds}: it lends them cash against their long positions,
lends them securities to sell short, holds and settles their portfolios, and
is paid a spread on every balance. \Cref{ch:m1:financing} takes it apart.
\end{definition}

\begin{definition}[Underwriting]\label{def:m1:the-sell-side:underwriting}
In \emph{underwriting}\index{underwriting} a bank buys newly issued shares or
bonds from the issuer and resells them to investors, earning a percentage of
the amount raised and bearing, for a few hours or days, the risk that the
issue does not sell. With \emph{advisory} work on mergers it forms
\emph{investment banking}, housed next to the markets division because the
new issues feed the trading desks.
\end{definition}

\begin{dated}{2026-09}{One markets division in its own numbers}\label{dat:m1:the-sell-side:gs}
Goldman Sachs reported 2025 net revenues of \$41.5 billion in Global Banking
\& Markets, out of \$58.3 billion for the firm. Within it: fixed income,
currencies and commodities \$14.5 billion, of which \$10.3 billion from
intermediation (making prices) and \$4.3 billion from financing; equities
\$16.5 billion, of which \$9.3 billion intermediation and \$7.2 billion
financing; investment banking fees \$9.3 billion (advisory \$4.7 billion,
equity \omterm{def:m1:the-sell-side:underwriting}{underwriting} \$1.8 billion, debt \omterm{def:m1:the-sell-side:underwriting}{underwriting} \$2.8 billion).
\end{dated}

\begin{omfigure}
\begin{tikzpicture}
\begin{axis}[width=11.5cm,height=5.2cm,xbar,bar width=7pt,
  xlabel={2025 net revenues (\$ billion)},
  ytick=data,yticklabels from table={figdata/markets-1/02-the-sell-side/gs_gbm_2025.csv}{label},
  yticklabel style={font=\small},tick label style={font=\small},label style={font=\small},
  y dir=reverse,xmin=0,xmax=12,enlarge y limits=0.1,grid=major,grid style={gray!25},
  table/col sep=comma]
\addplot[fill=omDef!60,draw=omDef] table[x=usd_bn,y=k]{figdata/markets-1/02-the-sell-side/gs_gbm_2025.csv};
\end{axis}
\end{tikzpicture}

\omcaption{Where one large markets division earned its revenue in 2025.
Financing --- lending against positions --- is more than a third of the
trading revenue. Data: the firm's earnings release, as in
\cref{dat:m1:the-sell-side:gs}.}
\end{omfigure}

\section{What regulation did to the dealer}

Before the 2008 crisis large dealers also traded for their own account on
dedicated proprietary desks. Two sets of
rules ended that and reshaped what remained.

\begin{definition}[Risk-weighted assets]\label{def:m1:the-sell-side:rwa}
A bank's \emph{risk-weighted assets}\index{risk-weighted assets} (RWA) are
its exposures, each multiplied by a regulatory weight reflecting its risk;
the bank must hold equity capital of at least a fixed percentage of their
sum. For a trading desk, the market-risk part of RWA is computed from the
riskiness of its positions under the Basel Committee's market-risk standard.
\end{definition}

\begin{method}[Judging a desk by its return on capital]\label{met:m1:the-sell-side:roc}
\begin{enumerate}
\item Take the desk's revenue less its direct costs.
\item Divide by the equity its positions consume: RWA times the bank's
  target capital ratio.
\item Compare with the bank's cost of equity, around ten percent.
\end{enumerate}
A position that ties up capital for weeks must earn far more than one that
turns over in minutes: this is why bank desks prefer flow and financing to
large, slow risk.
\end{method}

\begin{example}[Two desks]\label{ex:m1:the-sell-side:twodesks}
Desk A earns \$60 million on RWA of \$2 billion; desk B earns \$25 million on
RWA of \$400 million. At a 12\% capital ratio A uses \$240 million of equity
and returns 25\%; B uses \$48 million and returns 52\%. The bank grows B.
\end{example}

The second rule is a prohibition. Section~13 of the US Bank Holding Company
Act, added by the Dodd--Frank Act of 2010 and known as the \emph{Volcker
rule}, forbids banking entities to engage as principal in short-term
proprietary trading. Its exemptions define the modern dealer: \omterm{def:m1:the-sell-side:underwriting}{underwriting},
\emph{market making-related} activity, risk-mitigating hedging, trading on
behalf of customers, and trading in government obligations remain allowed.
A desk must therefore be able to show that its inventory is sized to
the ``reasonably expected near term demands of clients, customers, or
counterparties'' --- that it is a shop, not a fund. Much of the proprietary talent that left the banks after 2010
founded or joined the firms of \cref{ch:m1:what-a-trading-firm-does}.

\section{Tutorial: pricing the risk bid}\label{tut:m1:the-sell-side:riskbid}

\begin{tutorial}
\textbf{Goal.} Price the pension fund's block and check the $T/3$ rule by
simulation. \textbf{End state:} the histogram below, with 5\% of its mass
left of zero.
\begin{tutsteps}
\item \textbf{Code the two propositions.}
\omcode{code/markets-1/02-the-sell-side/python/blockbid.py}{29}{45}{Unwind time, impact, risk and the break-even discount.}\label{lst:m1:the-sell-side:formulas}
\item \textbf{Simulate the unwind.} Two hundred equal slices over $T$ days,
each sold at the then price; the dealer's result is the discount, less
impact, plus the average return of the slices.
\omcode{code/markets-1/02-the-sell-side/python/blockbid.py}{47}{54}{Monte Carlo of a linear unwind: the average of a random walk.}\label{lst:m1:the-sell-side:mc}
\item \textbf{Run and compare.} For the example's block the simulated
standard deviation is $1.64\%$ against a predicted $0.02\sqrt{2/3} =
1.63\%$, and $4.9\%$ of 20\,000 paths lose money at the 95\% discount.
\item \textbf{Read the histogram.} The mean, $2.7\%$, is the dealer's
expected profit; it is what the client pays for certainty.
\end{tutsteps}
\textbf{What to change next.} Raise \texttt{participation} to 0.25: the
discount falls to $2.3\%$, and the model says faster is always better. Why
does no desk sell at 50\% of volume? (Exercise~\ref{exo:m1:the-sell-side:8}.)
Then hedge: replace $\sigma$ by the stock's residual volatility.
\end{tutorial}

\begin{omfigure}
\begin{tikzpicture}
\begin{axis}[width=11cm,height=4.8cm,ybar,bar width=6.5pt,
  xlabel={dealer's profit (\% of block value)},ylabel={paths out of 20\,000},
  tick label style={font=\small},label style={font=\small},
  xmin=-4.2,xmax=9.2,ymin=0,table/col sep=comma]
\addplot[fill=omDef!40,draw=omDef] table[x=centre,y=count]{figdata/markets-1/02-the-sell-side/pnl_hist.csv};
\addplot[omThm,dashed,sharp plot,no markers] coordinates {(0,0) (0,2500)};
\end{axis}
\end{tikzpicture}

\omcaption{Profit of the risk bid at the 95\% break-even discount. The dashed
line is zero: one path in twenty falls to its left. Data: the tutorial's
simulation, seed 2.}
\end{omfigure}

\section{Build: the block pricer}\label{bld:m1:the-sell-side:blockbid}

\begin{build}
\bfield{Purpose} A pricing service the miniature firm's bank-style desk
calls when a client asks for a risk bid.
\bfield{Interface} \texttt{quote(shares, price, adv, sigma, participation,
confidence, hedge\_ratio=0.0)} returning the unwind time, impact, risk,
discount and bid price. \texttt{hedge\_ratio} is the fraction of variance
removed by an index hedge: the risk term uses $\sigma\sqrt{1-h}$.
\bfield{Rules} Reject a participation above 30\% and a block above five
days' volume (\texttt{ValueError}): outside that range the square-root rule
is not to be trusted. Round the bid \emph{down} to the cent.
\bfield{Acceptance tests} \texttt{code/firm/blockbid/tests/}: reproduces
\cref{ex:m1:the-sell-side:block} to the cent; a fully hedged block is priced
at impact only; limits raise.
\bfield{Stretch} Post the desk's realised profit on a simulated unwind to
the ledger of \cref{ch:m1:what-a-trading-firm-does} under \texttt{position}.
\end{build}

\begin{omsources}
\item Goldman Sachs, \emph{Full Year and Fourth Quarter 2025 Earnings
Results}, Form 8-K, January 2026.
\item Federal Deposit Insurance Corporation, ``Volcker Rule'' (overview of
Section 13 of the Bank Holding Company Act and its implementing rules);
\emph{Federal Register}, 31 July 2020, final rule on covered funds.
\item Basel Committee on Banking Supervision, \emph{Minimum capital
requirements for market risk}, January 2019.
\item L.~Harris, \emph{Trading and Exchanges}, Oxford University Press, 2003,
chapters 15 and 16 (block trading, dealers).
\item J.-P.~Bouchaud, J.~Bonart, J.~Donier and M.~Gould, \emph{Trades, Quotes
and Prices}, Cambridge University Press, 2018, chapter 12 (the square-root
law).
\end{omsources}

\section{Exercises}

\begin{exercise}[$\star$]\label{exo:m1:the-sell-side:1}
A dealer holds a \$30 million position in a stock with daily volatility
1.5\% and sells it linearly over three days. Give the standard deviation of
its result in dollars.
\end{exercise}

\begin{exercise}[$\star$]\label{exo:m1:the-sell-side:2}
A block of 500\,000 shares in a stock trading 5 million a day is unwound at
10\% participation. How long does the unwind take, and what is the impact
cost in percent if $\sigma = 1.5\%$ and $\kappa = 0.7$?
\end{exercise}

\begin{exercise}[$\star$]\label{exo:m1:the-sell-side:3}
An issuer raises \$800 million of new shares with an \omterm{def:m1:the-sell-side:underwriting}{underwriting} fee of
3.5\%. What do the banks earn? The shares are placed at \$40; on the first
day they close at \$46. How much did the issuer ``leave on the table'', and
who received it?
\end{exercise}

\begin{exercise}[$\star\star$]\label{exo:m1:the-sell-side:4}
Compute the 95\% break-even discount, in percent and in dollars a share, for
a block of 3 million shares at \$20 in a stock with $V = 6$ million and
$\sigma = 2.5\%$, at 10\% participation.
\end{exercise}

\begin{exercise}[$\star\star$]\label{exo:m1:the-sell-side:5}
Show that at fixed participation the break-even discount is proportional to
$\sqrt{Q}$. A dealer quotes 2\% for a block; what should it quote, other
things equal, for a block nine times larger?
\end{exercise}

\begin{exercise}[$\star\star$]\label{exo:m1:the-sell-side:6}
With the figures of \cref{dat:m1:the-sell-side:gs}: (a) what share of the
division's revenue came from financing, from intermediation and from
investment banking fees? (b) Which of the five revenue sources of
\cref{met:m1:what-a-trading-firm-does:sources} is each?
\end{exercise}

\begin{exercise}[$\star\star\star$]\label{exo:m1:the-sell-side:7}
\emph{Coding.} An index hedge removes 40\% of the variance of the stock of
\cref{ex:m1:the-sell-side:block}. (a) Compute the new break-even discount.
(b) The hedge costs 2 basis points of the block's value in futures
commissions and spread. Is it worth it?
\end{exercise}

\begin{exercise}[$\star\star\star$]\label{exo:m1:the-sell-side:8}
\emph{Find the flaw.} A junior trader observes that in
\cref{prop:m1:the-sell-side:discount} the discount falls as participation
rises, and proposes to unwind every block at 50\% of volume. Give two
reasons, one inside the model's assumptions and one outside them, why the
desk's discount would not fall as the formula says.
\end{exercise}

\section{Problem: The Risk Bid}

\begin{problem}[{Weekend problem --- forty seconds to price a block}]\label{pb:m1:the-sell-side:1}
You run the cash equities risk book at a bank. At 11:40 a pension fund asks
for a bid on 4 million shares of a stock trading at \$25, with a daily volume
of 8 million shares and a daily volatility of 1.8\%. Use $\kappa = 0.7$.

\medskip\noindent\textbf{Part I --- The mechanical price.}
\begin{enumerate}
\item What is the value of the block, and what fraction of daily volume?
\item At 10\% participation, how many days does the unwind last?
\item Compute the impact cost in percent and in dollars.
\item Compute the standard deviation of the unwind result in percent and in
dollars.
\item Compute the 95\% break-even discount and the corresponding bid price,
rounded down to the cent.
\end{enumerate}

\medskip\noindent\textbf{Part II --- Improving it.}
\begin{enumerate}[resume]
\item You are willing to be 20\% of the volume. Recompute the unwind time
and the discount.
\item An index hedge removes half the variance. Recompute the discount at
20\% participation.
\item What is your expected profit in dollars if you win the block at that
discount?
\item A competitor is known to bid at one standard deviation ($z = 1$)
instead of 1.645. What discount does it quote with the same hedge and
participation, and how often does it lose money?
\end{enumerate}

\medskip\noindent\textbf{Part III --- The client.}
\begin{enumerate}[resume]
\item The alternative for the fund is an agency execution over the same
$T$ days at a commission of 1 cent a share. Give its expected cost in
percent (impact plus commission) and the standard deviation of its outcome.
\item By how much does the risk bid of question 7 exceed that expected
cost? What does the fund buy with the difference?
\item The fund's trader is measured against the price at the time of the
decision. Explain why she may prefer the risk bid even at that premium.
\item Over the last year this client's blocks were followed, on average, by
a fall of 0.4\% in the stock relative to the index over the next two days.
What should you add to your discount, and what is this cost called?
\end{enumerate}

\medskip\noindent\textbf{Part IV --- The bank.}
\begin{enumerate}[resume]
\item The position adds \$150 million of RWA for three days. At a 12\%
capital ratio and a 10\% cost of equity, what does that capital cost in
dollars (use 252 days)?
\item Compare with the expected profit of question 8. Is capital the
binding cost of this trade?
\item The same client paid the bank \$3 million last year in financing and
derivatives. How low could you rationally bid?
\item You win the block and sell \$100 million of index futures against it.
In the first hour the index falls 1\% and the stock 3\%. Estimate the
mark-to-market loss on the pair before any sale.
\item Under the Volcker rule, what would you need to show an examiner about
this position?
\item State the \emph{named result}: the discount, in cents a share, at
which this desk breaks even with 95\% confidence after hedging and the
adverse-selection charge of question 13.
\item In one sentence, why is a risk bid an option the client owns?
\end{enumerate}
\end{problem}

\section{Interview questions}

\begin{interviewq}[$\star$ \iqroles{bank, trader}]\label{iq:m1:the-sell-side:1}
What is the difference between a \omterm{def:m1:the-sell-side:sales-trader}{sales-trader} and a trader? Who is
responsible for the P\&L of a client order in each case?
\end{interviewq}

\begin{interviewq}[$\star\star$ \iqroles{trader, researcher, bank}]\label{iq:m1:the-sell-side:2}
You must sell a position evenly over nine days. Roughly how much risk are
you carrying compared with holding it unchanged for the same nine days?
\end{interviewq}

\begin{interviewq}[$\star\star$ \iqroles{trader, bank}]\label{iq:m1:the-sell-side:3}
A client asks for a two-way price in a block of a stock you know nothing
about. What do you ask yourself in the ten seconds you have?
\end{interviewq}

\begin{interviewq}[$\star\star$ \iqroles{bank, researcher}]\label{iq:m1:the-sell-side:4}
Why did financing grow to more than a third of the trading revenue of large
banks after 2010, while their proprietary risk shrank?
\end{interviewq}

\begin{interviewq}[$\star\star\star$ \iqroles{researcher, trader}]\label{iq:m1:the-sell-side:5}
The cost of a block grows like the square root of its size. A client can
give you one block of 4 million shares or four blocks of 1 million on four
consecutive days. Which is cheaper for the client, and what assumption
decides it?
\end{interviewq}

\begin{interviewq}[$\star\star\star$ \iqroles{bank, researcher}]\label{iq:m1:the-sell-side:6}
Derive the variance of the time-average of a Brownian motion over $[0,T]$.
\end{interviewq}
